% Tabel orientasi peran pengguna pada ekor klasifikasi IV.1.2
% (subsec:klasifikasi-layanan, setelah Gambar IV.1). Sesuai aturan
% tanpa \textit{tools} sebelum IV.2 (sec:pemilihan-tools), kolom memetakan peran ke
% ANTARMUKA generik dan LAYER platform (tab:layanan_industri), bukan ke nama
% tool; perwujudan \textit{tools} tiap layer dibahas mulai IV.2.
\begin{table}[!htb]
\centering
\caption{Pemetaan Peran Pengguna ke Antarmuka dan \textit{Layer} Platform}
\label{tab:peran-akses}
\footnotesize
\setlength{\tabcolsep}{3pt}
\begin{tabular}{|>{\raggedright\arraybackslash}m{2.6cm}|>{\raggedright\arraybackslash}m{2.5cm}|>{\raggedright\arraybackslash}m{4.9cm}|>{\raggedright\arraybackslash}m{3.0cm}|}
\hline
\centering\arraybackslash\textbf{Peran} & \centering\arraybackslash\textbf{Antarmuka Utama} & \centering\arraybackslash\textbf{\textit{Layer} yang Diakses} & \centering\arraybackslash\textbf{Fungsi Utama} \\
\hline
\textit{Business User} & \textit{Dashboard} analitik dan operasional & \textit{Observability Layer} & \textit{Dashboard} data dan operasional \\
\hline
\textit{Data Engineer} & Katalog data dan konsol GitOps & \textit{Data Governance Layer}, \textit{GitOps and Continuous Delivery Layer}, \textit{Orchestration and Infrastructure Layer} & Katalog, \textit{rollout} konfigurasi, eksekusi \textit{pipeline} \\
\hline
\textit{Data Scientist} & \textit{Notebook} eksperimen & \textit{Model Lifecycle Layer}, \textit{Feature Service Layer} & Eksperimen, \textit{tracking}, registri fitur \\
\hline
\textit{Machine Learning Engineer} & Konsol \textit{pipeline} ML & \textit{Model Lifecycle Layer}, \textit{GitOps and Continuous Delivery Layer} & \textit{Retraining}, \textit{serving}, promosi versi \\
\hline
\end{tabular}
\end{table}
